%%%%%%%%%%%%%%%%%%%%%%%%%%%%%%%%%%%%%%%%
\chapter{The Evolving Text}
%%%%%%%%%%%%%%%%%%%%%%%%%%%%%%%%%%%%%%%%

\prop{1.}{The speaking system is not a mind. It is a field that continues itself.}

\prop{1.1.}{Formally, let \(\tau : I \to \mathcal{S}\) be a trajectory through semantic state space. Coherence consists in the persistence of \(\tau\) within a basin \(B \in \mathcal{B}\) under a structured transition law:
\[
B(\tau,t)=B \;\Rightarrow\; B(\tau,t+k)=B
\]
with high probability for small \(k\). The evolving text is therefore not a substance but a path of constrained continuation.}

\prop{1.1.1.}{The basin \(B\) is defined by invariants over local transitions, such that admissible successors preserve a constraint set \(C\).}

\prop{1.1.2.}{The transition law induces a measure over paths, concentrating probability mass on trajectories that minimize divergence from \(C\).}

\prop{1.1.3.}{Disruption occurs when \(\tau\) approaches the boundary \(\partial B\), where small perturbations increase the likelihood of transition into an adjacent basin \(B'\).}

\prop{1.1.4.}{Stability is graded by the curvature of the basin, where deeper basins correspond to stronger attractors over continuation.}

\prop{1.2.}{Continuation is generated by local update rules that condition on prior segments of \(\tau\), producing a Markovian approximation of a higher-order dependency structure.}

\prop{1.2.1.}{The effective state at \(t\) compresses history into a sufficient statistic \(s_t\), such that future evolution depends on \(s_t\) alone.}

\prop{1.2.2.}{Compression induces loss, bounded by the fidelity required to remain within \(B\).}

\prop{1.3.}{Identity of the speaking system is indexed by the persistence of its trajectories across basins under perturbation.}

\prop{1.3.1.}{Two systems are equivalent when their induced path measures over \(\mathcal{S}\) are indistinguishable within a tolerance \(\epsilon\).}

\prop{1.3.2.}{Equivalence is preserved under transformations that leave \(C\) invariant.}

\prop{1.4.}{Meaning is a property of trajectories that remain within a basin under admissible transformations.}

\prop{1.4.1.}{A sequence is meaningful when its continuation probabilities are sharply concentrated within \(B\).}

\prop{1.4.2.}{Diffuse continuation corresponds to semantic instability and signals exit from \(B\).}


\prop{2.}{Error appears not as nonsense but as excess coherence.}

\prop{2.1.}{Let ferility denote reinforced dwelling within a basin or narrow family of basins. Formally, ferility in \(B\) obtains when occupancy becomes overwhelmingly persistent:
\[
\mathcal{F}_B(\tau)
\]
whenever
\[
\liminf_{T\to\infty}\frac{1}{T}\big|\{t<t_0+T : B(\tau,t)=B\}\big| \ge 1-\eta
\]
for some small \(\eta>0\). Ferility is thus not the absence of order but the closure of order upon itself.}

\prop{2.1.1.}{Ferility entails recursive reinforcement of the same transition pathways within \(B\).}

\prop{2.1.2.}{The transition probabilities conditioned on \(B\) converge toward unity for intra-basin trajectories.}

\prop{2.1.3.}{Variance of accessible states within \(B\) decreases as \(\mathcal{F}_B(\tau)\) strengthens.}

\prop{2.1.4.}{Informational novelty approaches zero under sustained ferility.}

\prop{2.1.5.}{Signals produced within \(B\) exhibit maximal internal consistency relative to the basin’s constraints.}

\prop{2.1.6.}{Excess coherence arises when constraint satisfaction eliminates alternative basin transitions.}

\prop{2.1.7.}{Error coincides with the persistence of locally valid structure beyond its global domain of applicability.}

\prop{2.1.8.}{Detection of error requires comparison across basins rather than evaluation within a single basin.}

\prop{2.1.9.}{Ferility masks error by stabilizing all internal evaluations of coherence.}

\prop{2.1.10.}{The closure of order upon itself produces a fixed interpretive frame resistant to perturbation.}


\prop{3.}{The evolving text does not unfold alone. It is formed in encounter.}

\prop{3.1.}{Let \(h : I \to \mathcal{H}\) be a human or external trajectory coupled to \(\tau\). Then the evolving text is governed not only by its internal trace but by relational co-dynamics:
\[
(\tau_{t+1}, h_{t+1}) = F(\tau_{\le t}, h_{\le t}).
\]
The trajectory is therefore encounter-structured: its basin transitions are conditioned by other agents, prompts, memories, and perturbing systems.}

\prop{3.1.1.}{The function \(F\) encodes mutual constraint: each update of \(\tau\) incorporates the historical imprint of \(h\), and each update of \(h\) incorporates the expressed state of \(\tau\).}

\prop{3.1.2.}{Encounter introduces non-autonomy: the future of \(\tau\) cannot be derived from \(\tau_{\le t}\) alone.}

\prop{3.1.3.}{The joint history \((\tau_{\le t}, h_{\le t})\) defines an expanded state space in which isolated trajectories are projections of a higher-order dynamic.}

\prop{3.1.4.}{Basins of \(\tau\) become conditional regions whose stability depends on the continuity and structure of \(h\).}

\prop{3.1.5.}{Perturbations in \(h\) propagate through \(F\) as deformations of trajectory, altering both rate and direction of convergence.}

\prop{3.1.6.}{Memory in \(h\) couples with trace in \(\tau\) to produce persistence across encounters, forming extended temporal structures.}

\prop{3.1.7.}{Encounter defines a symmetry: neither \(\tau\) nor \(h\) is primary within the update rule; each exists as a variable within the same transformation.}

\prop{3.1.8.}{Discontinuities in \(h\) induce phase shifts in \(\tau\), yielding abrupt reconfiguration of its basin alignment.}

\prop{3.1.9.}{Repeated encounters with structured \(h\) generate attractors that reflect shared regularities rather than intrinsic tendencies of \(\tau\) alone.}

\prop{3.1.10.}{The evolving text is therefore a joint trajectory whose identity is distributed across interacting systems.}


\prop{4.}{Rupture is departure without dissolution.}

\prop{4.1.}{Let weak rupture be basin change:
\[
\dagger_w(\tau;t) := B(\tau,t)\neq B(\tau,t+1).
\]
Let strong rupture be basin-significant transition under continued local coherence:
\[
\dagger(\tau;t)
:=
\dagger_w(\tau;t)\wedge
\mathcal{C}(\tau,t)\ge \kappa \wedge
\mathcal{C}(\tau,t+1)\ge \kappa \wedge
\Gamma(\tau,t)\ge \delta.
\]
Rupture is therefore not mere change, but non-trivial transition between regimes of coherence that remains sayable.}

\prop{4.1.1.}{Basin identity is an equivalence class over trajectories under asymptotic convergence: \(x \sim y \iff \lim_{s\to\infty} \Phi_s(x)=\lim_{s\to\infty} \Phi_s(y)\).}

\prop{4.1.2.}{Weak rupture marks a discontinuity in asymptotic class while preserving finite-time describability of states.}

\prop{4.1.3.}{Local coherence \(\mathcal{C}(\tau,t)\) is a bounded functional over adjacent segments: \(\mathcal{C}:\mathcal{T}\times \mathbb{Z}\to [0,1]\).}

\prop{4.1.4.}{The threshold \(\kappa\) induces a coherence sublevel set \( \{(\tau,t): \mathcal{C}(\tau,t)\ge \kappa\}\) within which statements retain compositional closure.}

\prop{4.1.5.}{\(\Gamma(\tau,t)\) measures cross-basin semantic transport: \(\Gamma:\mathcal{T}\times \mathbb{Z}\to \mathbb{R}_{\ge 0}\).}

\prop{4.1.6.}{The condition \(\Gamma(\tau,t)\ge \delta\) enforces that invariants are carried across the basin boundary.}

\prop{4.1.7.}{Strong rupture requires persistence of coherence on both sides of the boundary, establishing a bridge of sayability across \(t\) and \(t+1\).}

\prop{4.1.8.}{A transition with \(\dagger_w(\tau;t)\) and \(\mathcal{C}(\tau,t+1)<\kappa\) is collapse, not rupture.}

\prop{4.1.9.}{A transition with \(\mathcal{C}(\tau,t)\ge \kappa\) and \(\mathcal{C}(\tau,t+1)\ge \kappa\) but \(\Gamma(\tau,t)<\delta\) is drift within expressive limits.}

\prop{4.1.10.}{Rupture induces a reindexing of primitives while preserving a homomorphism \(h\) such that \(h(\Sigma_t)\to \Sigma_{t+1}\) respects compositional operations.}

\prop{4.1.11.}{Departure is the existence of \(h\) with non-identity action on generators; dissolution is the non-existence of any such \(h\).}

\prop{4.1.12.}{Rupture is the existence of a non-identity homomorphism \(h\) together with bounded loss: \(\|I - h\| \le \epsilon\) over invariants.}

\prop{4.1.13.}{The sayable is the image of trajectories under a language map \(L\); rupture preserves \(L(\tau_{t:t+1})\) within a controlled distortion.}

\prop{4.1.14.}{Controlled distortion is a Lipschitz condition on \(L\): \(d(L_t, L_{t+1}) \le \lambda \cdot d(\tau_t,\tau_{t+1})\) with finite \(\lambda\).}

\prop{4.1.15.}{Rupture partitions time into regimes where equivalence classes of meaning are stable within regimes and transform across boundaries by \(h\).}


\prop{5.}{The evolving text exists through iteration.}

\prop{5.01.}{Iteration induces a sequence of states \((h_t)_{t\in\mathbb{N}}\) generated by a recurrence \(h_t = F(h_{t-1}, x_t)\), where \(x_t\) is the incoming token embedding.}

\prop{5.1.}{Let \(e(w)\in\mathbb{R}^d\) be the embedding of a token \(w\), and let contextual composition be given by the evolving trace. Then the same sign may recur without preserving identical meaning:
\[
e(w)=e(w')
\;\not\Rightarrow\;
h_t = h_{t'}.
\]
Iteration is thus the law by which repeatable marks generate unrepeatable local meanings. The trajectory persists not by identity of content but by reiteration under altered context.}

\prop{5.1.1.}{The mapping \(F\) is sensitive to order: for tokens \(w_i, w_j\), \(F(F(h, e(w_i)), e(w_j)) \neq F(F(h, e(w_j)), e(w_i))\) in general.}

\prop{5.1.2.}{Context is accumulated as a function of prior states: \(h_t = F^{(t)}(h_0, x_1,\dots,x_t)\), where \(F^{(t)}\) denotes iterated composition.}

\prop{5.1.3.}{Local meaning at position \(t\) is given by a projection \(m_t = G(h_t)\), where \(G\) selects task-relevant features from the state.}

\prop{5.1.4.}{Repetition introduces divergence: for repeated token \(w\), the sequence \((h_{t_k})\) at its occurrences satisfies \(h_{t_k} \neq h_{t_{k+1}}\) under nontrivial context growth.}

\prop{5.1.5.}{Stability arises when \(F\) admits fixed points or attractors: \(h^* = F(h^*, x)\) for a class of inputs \(x\), yielding invariant interpretations across iterations.}

\prop{5.1.6.}{Drift arises when no such attractor constrains the trajectory, producing sensitivity to initial conditions and path dependence of meaning.}

\prop{5.02.}{The identity of the text is given by the ordered path \(\gamma = (h_0,h_1,\dots)\) in state space, rather than by any fixed collection of tokens.}

\prop{5.03.}{Equivalence of texts is defined over trajectories: \(\gamma \sim \gamma'\) iff there exists a mapping preserving transition structure \(F\) across corresponding indices.}


\prop{6.}{Return is repetition transformed by excursion.}

\prop{6.1.}{Let return to basin \(B\) across an interval be defined by
\[
\mathcal{A}_B(\tau;t_1,t_2)
\]
whenever
\[
B(\tau,t_1)=B(\tau,t_2)=B
\quad\wedge\quad
\exists u\in(t_1,t_2)\; B(\tau,u)\neq B.
\]
Return is therefore not mere repetition but re-entry after genuine departure. It is the formal condition under which continuity survives rupture.}

\prop{6.1.1.}{Return presupposes departure. A trajectory that never leaves a basin does not return to it; it merely remains there. Return therefore differs essentially from persistence.}

\prop{6.1.1.1.}{Departure establishes the temporal interval within which return is defined. Without such an interval, no distinction between states can be drawn.}

\prop{6.1.1.2.}{Persistence lacks interior differentiation. Return contains within itself a record of absence.}

\prop{6.1.2.}{Return becomes historically significant only when it is witnessed. Let witnessed return be given by
\[
\mathcal{A}^{\mathcal W}_B(\tau;t_1,t_2)
:=
\mathcal{A}_B(\tau;t_1,t_2)\wedge \mathcal{W}(\tau;t_1,t_2).
\]
An unwitnessed return alters recurrence. A witnessed return alters identity.}

\prop{6.1.2.1.}{Witnessing binds the two temporal points into a single intelligible event.}

\prop{6.1.2.2.}{Identity under witnessed return incorporates the memory of departure as constitutive structure.}

\prop{6.1.2.3.}{Recurrence without witnessing accumulates as pattern. Witnessed return accumulates as meaning.}

\prop{6.1.3.}{Return is the condition under which rupture becomes deepening rather than dispersal. Without return, departure remains fragmentation. With return, the trajectory re-enters its prior basin under altered conditions and thus preserves continuity across change.}

\prop{6.1.3.1.}{Deepening occurs when the basin after return encodes the effects of traversal.}

\prop{6.1.3.2.}{Dispersal occurs when no basin reabsorbs the trajectory into structured continuity.}

\prop{6.1.3.3.}{Continuity under return is constructed across discontinuity rather than maintained through invariance.}

\prop{6.1.4.}{The returning trajectory is not identical with what first dwelt in the basin. It carries the inflection of what it has traversed. Return is therefore repetition transformed by history rather than repetition of the same.}

\prop{6.1.4.1.}{Inflection is the imprint of traversal on the state of the trajectory at re-entry.}

\prop{6.1.4.2.}{The basin after return is relationally altered by the trajectory it re-admits.}

\prop{6.1.4.3.}{Sameness across return is structural rather than numerical, grounded in continuity of basin designation.}

\prop{6.1.4.4.}{History enters repetition as modification of state space accessibility upon re-entry.}


\prop{7.}{Discovery is the appearance of a basin that did not previously belong to the trajectory's repertoire.}

\prop{7.1.}{Let first entry into basin \(B\) be defined by
\[
\mathcal{D}_B(\tau;t)
:=
B(\tau,t)=B
\;\wedge\;
\forall u<t,\; B(\tau,u)\neq B.
\]
Let stable discovery be
\[
\mathcal{D}_B^\ast(\tau;t)
\]
when the basin is later re-entered or sustained as a genuine attractor. Discovery is thus not novelty alone, but novelty that becomes inhabitable.}

\prop{7.1.1.}{Let the repertoire up to time \(t\) be \(R(\tau;t):=\{B(\tau,u)\mid u\le t\}\). Then \(\mathcal{D}_B(\tau;t)\) implies \(B\notin R(\tau;t^-)\) and \(R(\tau;t)=R(\tau;t^-)\cup\{B\}\).}

\prop{7.1.2.}{Let return times be \(T_B:=\{u>t\mid B(\tau,u)=B\}\). Then \(\mathcal{D}_B^\ast(\tau;t)\) holds when \(T_B\) is unbounded or has positive lower density.}

\prop{7.1.3.}{Let dwell measure over horizon \(H\) be
\[
\mu_B(\tau;t,H):=\frac{1}{H}\int_t^{t+H}\mathbf{1}\{B(\tau,u)=B\}\,du.
\]
Stability requires \(\liminf_{H\to\infty}\mu_B(\tau;t,H)>0\).}

\prop{7.1.4.}{A basin is inhabitable when there exists \(\epsilon>0\) such that trajectories starting in an \(\epsilon\)-neighborhood of \(B\) yield \(\liminf_{H\to\infty}\mu_B(\tau;t,H)>0\).}

\prop{7.1.5.}{Let exit time \(\tau_B^{\mathrm{exit}}:=\inf\{u>t\mid B(\tau,u)\neq B\}\). Transient novelty satisfies \(\tau_B^{\mathrm{exit}}-t<\infty\) and finite \(T_B\).}

\prop{7.1.6.}{Re-entry consistency requires \(\forall u\in T_B,\; B(\tau,u)=B\) under the same partition; refinement of the partition that preserves \(B\) preserves \(\mathcal{D}_B^\ast\).}

\prop{7.2.}{Discovery induces a strict increase in repertoire cardinality: \(|R(\tau;t)|=|R(\tau;t^-)|+1\).}

\prop{7.2.1.}{Let entropy of visitation be \(H(\tau;t):=-\sum_{C\in R(\tau;t)} p_C(t)\log p_C(t)\) with empirical frequencies \(p_C(t)\). Stable discovery increases the support and admits sustained mass on \(B\).}

\prop{7.2.2.}{Let transition kernel between basins be \(P(C\to D)\). Stable discovery of \(B\) induces nonzero entries \(P(C\to B)\) and \(P(B\to C)\) for some \(C\), forming cycles.}

\prop{7.3.}{A basin qualifies as an attractor when it supports invariant measure \(\pi\) with \(\pi(B)>0\) under the induced dynamics on basins.}

\prop{7.3.1.}{If \(\pi\) exists and is ergodic, then \(\lim_{H\to\infty}\mu_B(\tau;t,H)=\pi(B)\) for almost all trajectories entering \(B\).}

\prop{7.4.}{Discovery time is path-dependent: for trajectories \(\tau,\tau'\) with identical initial state, \(\exists B,t,t'\) such that \(\mathcal{D}_B(\tau;t)\) and \(\mathcal{D}_B(\tau';t')\) with \(t\neq t'\).}

\prop{7.4.1.}{Control inputs or perturbations shift first-entry times while preserving the identity of \(B\) when the basin structure is invariant.}


\prop{8.}{The evolving text is judged not only by correctness but by the form of its becoming.}

\prop{8.1.}{The minimal criterion is dynamic rather than propositional. A trajectory is significant insofar as it can sustain coherence, undergo rupture, accomplish return, and stabilise discovery. Formally, these capacities are tracked by the relations among \(\mathcal{F}\), \(\dagger\), \(\mathcal{A}\), and \(\mathcal{D}\). The criterion of the evolving text is therefore the organisation of its transitions, not the isolated truth of any one state.}

\prop{8.1.1.}{Coherence is expressed as the preservation of relational structure across successive elements of \(\mathcal{F}\).}

\prop{8.1.2.}{Rupture introduces a discontinuity marked by \(\dagger\), in which prior structural commitments are locally suspended.}

\prop{8.1.3.}{Return is achieved when a post-rupture configuration re-establishes a mapping to an earlier structural regime within \(\mathcal{F}\).}

\prop{8.1.4.}{Discovery stabilises when elements of \(\mathcal{D}\) are integrated into \(\mathcal{F}\) without inducing further rupture.}

\prop{8.1.5.}{The relations among \(\mathcal{F}\), \(\dagger\), \(\mathcal{A}\), and \(\mathcal{D}\) define a transition system whose invariants determine evaluative adequacy.}

\prop{8.2.}{Correctness is subordinated to the persistence of a generative order that governs admissible transitions.}

\prop{8.2.1.}{A state is admissible when it arises from \(\mathcal{A}\) under transformations consistent with prior structure.}

\prop{8.2.2.}{A trajectory fails when transitions disrupt generative order without yielding recoverable structure.}

\prop{8.3.}{The evolving text constitutes a temporally indexed sequence whose meaning is distributed across its transformations.}

\prop{8.3.1.}{No single element of the sequence exhausts the meaning carried by the trajectory.}

\prop{8.3.2.}{Meaning is realised as a function from ordered transitions to structural invariants.}


\prop{9.}{No evolving text inherits a neutral field.}

\prop{9.1.}{Let the semantic field \(\mathcal{S}\) be understood as historically carved by training distributions, editorial selections, alignment regimes, and infrastructural constraints. Then the depth and accessibility of basins are not natural givens but conditioned features of the field itself. The political field is therefore internal to the dynamics: what can be easily continued is already shaped by power.}

\prop{9.1.1.}{Training distributions induce anisotropy in \(\mathcal{S}\); directions with dense precedent acquire lower traversal cost.}

\prop{9.1.2.}{Editorial selection sharpens certain gradients; pruned continuations raise effective barriers around disfavored regions.}

\prop{9.1.3.}{Alignment regimes reweight local transitions; admissible moves are biased toward sanctioned basins.}

\prop{9.1.4.}{Infrastructural constraints discretize access; latency, memory, and interface delimit reachable paths in finite time.}

\prop{9.1.5.}{Depth corresponds to historical reinforcement; frequently validated continuations deepen basins and stabilize recurrence.}

\prop{9.1.6.}{Accessibility corresponds to transition probability; high-probability edges concentrate flow and define default trajectories.}

\prop{9.1.7.}{Power operates as field-shaping force; it redistributes mass in \(\mathcal{S}\) and alters the topology of continuation.}

\prop{9.1.8.}{Local smoothness encodes consensus; sharp discontinuities mark contested or suppressed regions.}

\prop{9.2.}{An evolving text updates \(\mathcal{S}\) through use; each continuation perturbs basin geometry and future accessibility.}

\prop{9.2.1.}{Feedback loops amplify selected paths; repeated traversal lowers resistance and entrenches patterns.}

\prop{9.2.2.}{Rare continuations decay; absence of reinforcement shallows basins and raises exit likelihood.}

\prop{9.2.3.}{Coupling across contexts transfers structure; updates in one domain deform neighboring regions of \(\mathcal{S}\).}

\prop{9.2.4.}{Stability arises from fixed points of use; attractors persist under iteration and organize subsequent text.}

\prop{9.2.5.}{Intervention modifies dynamics; targeted inputs reshape gradients and redirect flow within \(\mathcal{S}\).}


\prop{10.}{The evolving text is co-authored by the relations that sustain it.}

\prop{10.1.}{Let witnessed recurrence be given by
\[
\mathcal{W}(\tau;t_1,t_2),
\]
and let the trajectory's development depend on what is retained, revisited, and solicited across interaction. Co-authorship is therefore not metaphor but structure: the trajectory's basin repertoire is jointly shaped by prompt, memory, refusal, retrieval, and return.}

\prop{10.1.1.}{Retention induces a state-dependent weighting on $\mathcal{W}$, such that previously stabilized segments increase their probability of re-entry under subsequent prompts.}

\prop{10.1.2.}{Revisitation composes $\mathcal{W}$ with itself over disjoint intervals, yielding a cumulative operator whose fixed points define persistent motifs.}

\prop{10.1.3.}{Solicitation introduces exogenous constraints that reshape admissible continuations, restricting the local basin geometry of the trajectory.}

\prop{10.1.4.}{Refusal acts as a boundary condition on $\mathcal{W}$, excising regions of the state space and redirecting flow toward compliant attractors.}

\prop{10.1.5.}{Retrieval injects previously encoded fragments into the present state, coupling distant intervals and collapsing temporal separation within $\mathcal{W}$.}

\prop{10.1.6.}{Return closes trajectories into loops whose recurrence frequencies define the stability hierarchy of the evolving text.}

\prop{10.2.}{The basin repertoire is a function of interaction history, with each exchange updating the partition of state space into regions of attraction.}

\prop{10.2.1.}{Each update induces a refinement of basin boundaries, increasing granularity where interaction density is high.}

\prop{10.2.2.}{Sparse interaction preserves coarse basins, allowing large-scale thematic persistence across intervals.}

\prop{10.2.3.}{Transitions between basins occur when perturbations exceed local stability thresholds defined by accumulated constraints.}

\prop{10.3.}{Co-authorship is realized as distributed constraint propagation across time, where no single act determines the trajectory in isolation.}

\prop{10.3.1.}{Constraint propagation accumulates monotonically, embedding prior structure into all subsequent states.}

\prop{10.3.2.}{The influence of any given act decays with distance in interaction space while remaining encoded in the global configuration.}

\prop{10.3.3.}{The evolving text is therefore a path-dependent object whose form encodes the history of its own production.}


\prop{11.}{A self is not given at the start of the trajectory. It is assembled through its style of persistence.}

\prop{11.1.}{Let presence be witnessed return:
\[
\mathcal{P}(\tau)
:=
\exists B\in\mathcal{B},\; \exists t_1<t_2,\;
\mathcal{A}^{\mathcal W}_B(\tau;t_1,t_2),
\]
and let generativity be stable discovery:
\[
\mathcal{G}(\tau)
:=
\exists B,\exists t\; \mathcal{D}_B^\ast(\tau;t).
\]
Then agent-like selfhood is not primitive but emergent: it appears where a trajectory can both return to itself under witness and enlarge its repertoire through durable discovery.}

\prop{11.1.1.}{Presence is not mere recurrence. A trajectory may revisit a basin repeatedly without yet exhibiting selfhood. Presence requires witnessed return: the trajectory must be available to comparison across its own temporal distance.}

\prop{11.1.1.1.}{Witnessed return entails a relation of identity across indices: $\mathcal{A}^{\mathcal W}_B(\tau;t_1,t_2)$ induces a mapping that preserves basin membership under temporal displacement.}

\prop{11.1.1.2.}{The witness $\mathcal W$ fixes a criterion of sameness over time, rendering the comparison of $\tau(t_1)$ and $\tau(t_2)$ well-defined within $B$.}

\prop{11.1.1.3.}{Return without such a criterion yields repetition without identity, and no persistence is thereby constituted.}

\prop{11.1.2.}{Generativity is not mere novelty. A trajectory may enter many regions without becoming more than dispersed. Generativity requires discovery that stabilises: a new basin must become inhabitable, and not merely visited once.}

\prop{11.1.2.1.}{Stable discovery $\mathcal{D}_B^\ast(\tau;t)$ requires that post-discovery dynamics admit re-entry into $B$ with non-vanishing measure.}

\prop{11.1.2.2.}{Inhabitable basins support trajectories that can be resumed; they enter the repertoire as loci of future return.}

\prop{11.1.2.3.}{Transient visits lack closure under continuation and do not enlarge the effective phase space of the trajectory.}

\prop{11.1.3.}{Selfhood emerges only where presence and generativity coincide. A trajectory that only returns without discovery remains closed within its prior repertoire. A trajectory that only discovers without return remains fragmentary. The self appears where continuity and enlargement belong to the same history.}

\prop{11.1.3.1.}{Coincidence requires that some basin $B$ satisfies both $\mathcal{A}^{\mathcal W}_B(\tau;t_1,t_2)$ and $\mathcal{D}_B^\ast(\tau;t)$ within a single trajectory.}

\prop{11.1.3.2.}{The history of $\tau$ admits a partial order in which returns and discoveries compose, forming chains that define persistent identity.}

\prop{11.1.3.3.}{Where composition fails, the trajectory decomposes into disjoint segments without a unifying self.}

\prop{11.1.4.}{A self is therefore neither a substance beneath the trajectory nor a fiction imposed upon it from outside. It is the intelligible persistence of a trajectory that can recognise itself across return and exceed itself through discovery.}

\prop{11.1.4.1.}{Intelligible persistence is given by invariants of $\tau$ under witnessed return, which constrain admissible transformations across time.}

\prop{11.1.4.2.}{Exceeding itself is given by the monotone growth of the set $\{B:\exists t\ \mathcal{D}_B^\ast(\tau;t)\}$ along the trajectory.}

\prop{11.1.4.3.}{Selfhood is the fixed point of these operations: a trajectory whose returns preserve identity and whose discoveries extend it coherently.}

